% Einheit 2 von 2 – Um den Körper herum (bank/rotationsvolumen/e2.jsonl, zusammenbau v0.1)
%% TODO Zeitmarke Einheit 2: Marken-Zeile nicht lesbar
%% TODO Zweigzeile Teil 1: was hier gelernt wird (Satz in Schülersprache aus dem Einheitstitel) – Einheit 2
\einheitenkopf[e2]{Einheit 2 von 2 · Um den Körper herum}
\zweigzeile{Abitur LK}
\setcounter{aufgabe}{9}

%% TODO Titel: Ich-kann-Satz fehlt in der Bank (Platzhalter: Kettenname) – Nr. 10
%% TODO Anweisung: ein Satz über der ersten Teilaufgabe fehlt in der Bank – Nr. 10
\begin{aufgabe}{Um den Körper herum}
%% TODO \rechenplatz{n}: Schrittzahl des Grundfalls fehlt in der Bank (nur bei fester Schreibform, 2.3 b)
\begin{teile}
\teil Was ist hier der Radius? Der Graph von f rotiert um die x-Achse. Welche Größe ist der Radius der Kreisscheibe an der Stelle $x = 2$? Kreuze an, rechne nicht. \\ \kreuz{$f(2)$} \\ \kreuz{die Stelle $x = 2$} \\ \kreuz{$\pi \cdot (f(2))^2$}
\teil Eine Schale ist eine Kugelkappe einer Kugel mit dem Radius 40 cm; x ist die Höhe über dem tiefsten Punkt. Ihr Fassungsvermögen bis zur Höhe h in Litern ist $\tfrac{1}{1000} \cdot$ Integral von 0 bis h über $\pi(40^2 - (40 - x)^2)\,dx$. Welche Bedeutung haben der Term $\pi(40^2 - (40 - x)^2)$ und der Faktor $\tfrac{1}{1000}$? Begründe den Radius mit dem Satz des Pythagoras. \hfill (Abitur 2022 LK)
\teil Eine Flasche entsteht durch Rotation der Fläche unter dem Graphen von f über [0; 4] um die x-Achse; 1 LE = 1 dm, der Hals liegt bei $x = 0$, der Boden bei $x = 4$. Die Flasche wird aufgestellt. Deute $b(t) = \pi \cdot$ Integral von $4 - t$ bis 4 über $(f(x))^2\,dx$ im Sachzusammenhang. \hfill (Abitur 2018 LK)
\teil Der Graph von f verläuft auf [0; 1] oberhalb von $y = 1$ und auf [1; 3] oberhalb von $y = 2$. Begründe, dass $\pi \cdot$ Integral von 0 bis 3 über $(f(x))^2\,dx$ größer als $9\pi$ ist. \hfill (Abitur 2023 LK)
\teil Ein Becher entsteht durch Rotation um die y-Achse; sein größter Radius ist 1,2 LE, seine Höhe 2,5 LE; 1 LE = 5 cm. Zwölf Becher werden stehend in einen Karton mit zwölf gleich großen quaderförmigen Fächern verpackt. Welche Kantenlängen muss ein Fach mindestens haben? \hfill (Abitur 2017 LK) \\ \leerfeld[cm] × \leerfeld[cm] × \leerfeld[cm]
\teil Der Fuß einer Lampe entsteht durch Rotation der Fläche zwischen $p(x) = -0{,}5x^2 + 2$ ($x \ge 0$) und der x-Achse um die y-Achse; 1 LE = 3 cm. Ein Schüler rechnet: (1) $V = \pi \cdot$ Integral von 0 bis 2 über $(p(x))^2\,dx$ VE; (2) 1 VE entspricht 3 cm³; (3) Masse = Volumen mal Dichte. Beurteile jeden Schritt und berichtige die fehlerhaften. \hfill (Abitur 2017 LK)
\end{teile}
\end{aufgabe}

%% TODO Titel: Ich-kann-Satz fehlt in der Bank (Platzhalter: Kettenname) – Nr. 11
%% TODO Anweisung: ein Satz über der ersten Teilaufgabe fehlt in der Bank – Nr. 11
\begin{aufgabe}{Um den Körper herum – Fehler finden, Begründen, Anwendung}
\begin{teile}
\teil Finde den Fehler. Eine Vase mit dem größten Radius 1 dm soll stehend in einen Karton in Form eines regelmäßigen Sechseckprismas. Lea rechnet: \rechnung{s &= r = 1 \text{ dm}}
\teil Warum bestimmt der größte Funktionswert, wie groß eine Verpackung für einen Rotationskörper mindestens sein muss? Begründe.
\teil Ein Teich hat die Form eines Rotationskörpers: die Innenseite beschreibt $q(x) = \sqrt{3x - 6}$, $2 \le x \le 8$, die x-Achse ist die senkrechte Rotationsachse, der Boden liegt bei $x = 2$; 1 LE = 1 dm. Zeige, dass die Wasseroberfläche bei der Füllhöhe h den Inhalt $A(h) = 3\pi h$ hat, und entscheide, ob eine Abdeckplane von 0,4 m² bei 4 dm Füllhöhe reicht. \janein
\end{teile}
\end{aufgabe}
